% ============================================================
% qwmo_figure_final.tex
% Figure 1 — QWMO Operators
% Elsevier-compatible TikZ figure
% ============================================================

\begin{figure*}[t]
\centering

\begin{tikzpicture}[
    >=Stealth,
    thick,
    scale=0.86,
    every node/.style={scale=0.86},
    agent/.style={
        circle,
        draw=black,
        line width=0.45pt,
        inner sep=2.35pt
    },
    boxnote/.style={
        draw=gray!45,
        fill=gray!6,
        rounded corners,
        font=\scriptsize,
        align=left,
        inner sep=4pt
    }
]

% ============================================================
% PANEL SEPARATORS
% ============================================================

\draw[gray!35, dashed, line width=0.6pt]
    (5.8,-1.15) -- (5.8,5.0);

\draw[gray!35, dashed, line width=0.6pt]
    (11.6,-1.15) -- (11.6,5.0);

% ============================================================
% PANEL A — ADAPTIVE ORBITAL SAMPLING
% ============================================================

\node[font=\bfseries]
    at (2.9,4.7)
    {(a) Adaptive Orbital Sampling};

% Explorative agent
\node[
    agent,
    fill=red!55,
    label={[font=\scriptsize]below:
    Explorative ($q_i\!\approx\!0$)}
]
    (explore) at (1.0,1.5) {};

% Wide orbital
\draw[red!25, fill=red!5, opacity=0.65]
    (explore) circle (1.55);

\draw[red!45, dashed, fill=red!10, opacity=0.55]
    (explore) circle (0.9);

\draw[->, red!75]
    (explore) -- ++(0,1.55)
    node[right, font=\scriptsize]
    {$\sigma_{\max,i}$};

\draw[->, red!55, dashed]
    (explore) -- ++(0,0.9)
    node[left, font=\scriptsize]
    {$\sigma_i(t)$};

% Sample point from broad orbital
\node[
    agent,
    fill=red!35,
    inner sep=1.45pt,
    label={[font=\scriptsize]right:$\mathbf{x}'$}
]
    (sampleA) at (1.6,2.7) {};

\draw[->, red!55, dashed]
    (explore) -- (sampleA);

% Exploitative agent
\node[
    agent,
    fill=blue!60,
    label={[font=\scriptsize]below:
    Exploitative ($q_i\!\approx\!1$)}
]
    (exploit) at (4.5,1.5) {};

% Narrow orbital
\draw[blue!30, fill=blue!5, opacity=0.7]
    (exploit) circle (0.55);

\draw[->, blue!75]
    (exploit) -- ++(0,0.55)
    node[right, font=\scriptsize]
    {$\sigma_i(t)$};

% Sample point from narrow orbital
\node[
    agent,
    fill=blue!35,
    inner sep=1.45pt,
    label={[font=\scriptsize]right:$\mathbf{x}'$}
]
    (sampleB) at (4.7,1.95) {};

\draw[->, blue!55, dashed]
    (exploit) -- (sampleB);

% Formula box
\node[boxnote]
    at (2.9,-0.5)
{
$\sigma_{\max,i}=\gamma(u-\ell)(2-q_i)$\\
$\sigma_i(t)=
\max\!\left\{\sigma_{\max,i}e^{-c_i t/T_{\max}},
\sigma_{\lfloor}\right\}$\\
$c_i=c_{\mathrm{base}}q_i$
};

% ============================================================
% PANEL B — PAULI-INSPIRED EXCLUSION
% ============================================================

\node[font=\bfseries]
    at (8.7,4.7)
    {(b) Pauli-Inspired Exclusion};

% Better agent
\node[
    agent,
    fill=green!50!black,
    label={[font=\scriptsize]above:$\mathbf{x}_g$}
]
    (xg) at (8.0,2.5) {};

% Weaker agent
\node[
    agent,
    fill=orange!65,
    label={[font=\scriptsize]below:$\mathbf{x}_b$}
]
    (xb) at (7.2,1.0) {};

% Exclusion radius
\draw[gray!45, dashed]
    (xg) circle (1.8);

\node[font=\scriptsize, color=gray]
    at (9.6,3.8)
    {$\varepsilon$};

% Collision vector
\draw[->, black, line width=1.1pt]
    (xg) -- (xb)
    node[midway, left, font=\scriptsize]
    {$\mathbf{v}$};

% Orthogonal displacement
\draw[->, blue!75, line width=1.45pt]
    (xb) -- ++(1.2,-0.64)
    node[right, font=\scriptsize]
    {$\mathbf{p}\perp\mathbf{v}$};

% New position
\node[
    agent,
    fill=orange!40,
    inner sep=1.45pt,
    label={[font=\scriptsize]right:$\mathbf{x}_b'$}
]
    (xbnew) at (8.65,0.36) {};

\draw[
    ->,
    orange!70,
    dashed,
    line width=1.1pt
]
    (xb) -- (xbnew)
    node[midway, below right, font=\scriptsize]
    {$\alpha$};

% Formula box
\node[boxnote]
    at (8.7,-0.5)
{
$\mathbf{p}=
\mathbf{r}-
\dfrac{\mathbf{r}^{\top}\mathbf{v}}
{\|\mathbf{v}\|^2+10^{-10}}\mathbf{v}$\\
$\alpha\sim\mathcal{U}(0.5\varepsilon,1.5\varepsilon)$
};

% ============================================================
% PANEL C — ADAPTIVE QUANTUM ESCAPE
% ============================================================

\node[font=\bfseries]
    at (14.8,4.7)
    {(c) Adaptive Quantum Escape};

% Energy landscape
\draw[black, line width=1.25pt]
plot[smooth, tension=0.7] coordinates {
    (12.0,3.5)
    (12.5,2.8)
    (13.0,0.7)
    (13.5,2.2)
    (14.0,2.8)
    (14.6,1.8)
    (15.2,-0.3)
    (15.8,1.5)
    (16.5,3.2)
};

% Stagnating agent
\node[
    agent,
    fill=purple!65,
    label={[font=\scriptsize, color=purple]left:$k_i>k_s$}
]
    (trapped) at (13.0,0.7) {};

% Escape probability
\node[
    font=\scriptsize,
    color=purple!70
]
    at (13.1,1.55)
{
$P_{\mathrm{esc}}=e^{-2\kappa(t)L_{\mathrm{norm}}}$
};

% Global best
\node[
    agent,
    fill=black,
    inner sep=3pt,
    label={[font=\scriptsize]below:$\mathbf{x}^{*}$}
]
    (gopt) at (15.2,-0.3) {};

% Random exploration target
\node[
    agent,
    fill=gray!60,
    inner sep=2.2pt,
    label={[font=\scriptsize]right:$\mathbf{x}_r$}
]
    (randp) at (16.05,1.8) {};

% Wavy escape path
\draw[
    purple!50,
    line width=1pt,
    decoration={
        snake,
        amplitude=1.4pt,
        segment length=6pt
    },
    decorate
]
    (trapped) -- (14.55,0.7);

% Hybrid stochastic relocation arrows
\draw[
    ->,
    orange!85,
    line width=1.7pt,
    dashed,
    bend left=30
]
    (trapped) to
    node[above, font=\scriptsize, color=orange!80]
    {$\beta$}
    (gopt);

\draw[
    ->,
    gray!70,
    line width=1.4pt,
    dashed,
    bend left=12
]
    (trapped) to
    node[above, font=\scriptsize]
    {$1-\beta$}
    (randp);

% Formula box
\node[boxnote]
    at (14.8,-0.5)
{
$\kappa(t)=\kappa_0 t/T_{\max}$\\
$\mathbf{x}'=\beta(\mathbf{x}^{*}+\boldsymbol{\eta})
+(1-\beta)\mathbf{x}_r$
};

\end{tikzpicture}

\caption{
Schematic illustration of the three operators composing QWMO.
(a) Adaptive Orbital Sampling dynamically adjusts the Gaussian
sampling radius according to solution quality, balancing
exploration and exploitation.
(b) Pauli-Inspired Exclusion applies orthogonal displacement
when neighboring agents violate the exclusion radius
$\varepsilon$, promoting spatial diversity.
(c) Adaptive Quantum Escape enables stagnating agents
to leave local minima through probabilistic
$\beta$-hybrid stochastic relocation between
the global-best region and random exploration targets.
The terminology is used as an optimization analogy rather than
a direct physical simulation of quantum systems.
The combined interaction of the three operators establishes
the exploration--exploitation balance of QWMO.
}
\label{fig:qwmo_operators}

\end{figure*}